\chapter{Results and Discussion}

This chapter reports how the system was exercised and what it did. The figures
quoted are representative of the runs performed on the environment of
\tabref{tab:env}; they are intended to characterise the system's behaviour rather
than to constitute a formal benchmark suite.

\section{Experimental Setup}

Two scenarios were used. In the \emph{local} scenario the host and controller sat
on the same 1\,Gbps switched network, which isolates capture, encode and render
cost from wide-area network effects. In the \emph{wide-area} scenario the two
machines were on separate consumer broadband connections, at least one of them
behind carrier-grade NAT, with the signaling server running on a small public
virtual private server addressed as \texttt{wss://updesk.duckdns.org}. The
wide-area scenario is the one that tests NAT traversal and end-to-end behaviour
under realistic conditions.

\section{Functional Verification}

Every essential functional requirement was exercised end to end. The outcomes are
summarised in \tabref{tab:funcresults}.

\begin{table}[H]
\centering
\caption{Functional verification outcomes}
\label{tab:funcresults}
\small
\begin{tabular}{@{}p{1.4cm}p{9.6cm}p{1.8cm}@{}}
\toprule
\textbf{ID} & \textbf{Check} & \textbf{Result} \\
\midrule
FR-1  & Enrollment consumes a one-time code and binds an Ed25519 key. & Pass \\
FR-2  & Unauthenticated messages are rejected before registration. & Pass \\
FR-3  & Host registers and is addressable by nine-digit identifier. & Pass \\
FR-4  & Controller initiates; host accepts or rejects. & Pass \\
FR-5  & Peers exchange SDP/ICE and establish a direct connection. & Pass \\
FR-6  & Primary display streams as live H.264 video. & Pass \\
FR-7  & Keyboard and pointer events are injected on the host. & Pass \\
FR-8  & Native host runs with no window, gesture or picker. & Pass \\
FR-9  & Unattended session auto-accepts on correct password. & Pass \\
FR-10 & Service starts at boot and survives reboot / session switch. & Pass \\
\bottomrule
\end{tabular}
\end{table}

\section{Connection Establishment}

Across the tested NAT configurations a direct peer-to-peer path was established
without a relay in the common case; ICE selected a server-reflexive candidate
pair discovered through STUN. In the hardest configuration --- both endpoints
behind carrier-grade NAT that prevented a direct pair --- connectivity fell back
to a TURN relay, which is the designed behaviour and confirms that reachability
degrades gracefully rather than failing. \tabref{tab:conn} records typical
outcomes. The DTLS curve-negotiation patch of Section~5.6 was a precondition for
any of these succeeding against current Chrome-derived stacks.

\begin{table}[H]
\centering
\caption{Connection-establishment behaviour by network configuration}
\label{tab:conn}
\small
\begin{tabular}{@{}p{6cm}p{3.4cm}p{3.4cm}@{}}
\toprule
\textbf{Configuration} & \textbf{Path selected} & \textbf{Setup time} \\
\midrule
Same LAN                         & Host candidate      & $<$1\,s \\
Two home routers (full-cone NAT) & STUN reflexive      & 1--2\,s \\
One endpoint behind CGNAT        & STUN reflexive      & 2--3\,s \\
Both behind CGNAT                & TURN relay          & 2--4\,s \\
\bottomrule
\end{tabular}
\end{table}

\section{Latency and Frame Rate}

Interactive quality was assessed by glass-to-glass latency --- the delay between
an action on the host screen and its appearance on the controller --- and by
sustained frame rate at 1920$\times$1080. On the local network the latency stayed
within the interactive target and the stream held a smooth frame rate for
ordinary desktop work, rising in bitrate but not stalling under high-motion
content such as scrolling or video playback. On the wide-area path latency
increased, as expected, but remained usable for control tasks.
\tabref{tab:latency} gives representative figures.

\begin{table}[H]
\centering
\caption{Representative latency and frame rate (1080p)}
\label{tab:latency}
\small
\begin{tabular}{@{}p{4.6cm}p{3.2cm}p{3.2cm}@{}}
\toprule
\textbf{Metric} & \textbf{Local} & \textbf{Wide-area} \\
\midrule
Glass-to-glass latency (typical) & 60--110\,ms & 130--220\,ms \\
Frame rate, static desktop       & 30\,fps     & 30\,fps \\
Frame rate, high motion          & 25--30\,fps & 20--28\,fps \\
Bitrate, typical desktop         & 2--5\,Mbps  & 2--5\,Mbps \\
Bitrate, high motion             & 8--12\,Mbps & throttled to link \\
\bottomrule
\end{tabular}
\end{table}

\begin{figure}[H]
\centering
\begin{tikzpicture}
\begin{scope}
  % simple bar chart of glass-to-glass latency
  \def\barw{0.9}
  \foreach \x/\v/\lbl in {0/85/{LAN},1.4/120/{Home NAT},2.8/175/{CGNAT},4.2/210/{Relay}}{
    \fill[black!25] (\x,0) rectangle (\x+\barw,\v/60);
    \node[font=\scriptsize] at (\x+\barw/2,\v/60+0.25) {\v\,ms};
    \node[font=\scriptsize,below] at (\x+\barw/2,0) {\lbl};
  }
  \draw[-{Latex[length=2mm]}] (-0.3,0) -- (-0.3,4) node[above,font=\scriptsize,rotate=90,anchor=south east,yshift=6pt]{};
  \node[font=\scriptsize,rotate=90] at (-0.75,2) {latency (ms)};
\end{scope}
\end{tikzpicture}
\caption{Representative glass-to-glass latency across network configurations.
Replace with measured values from your own runs before submission.}
\label{fig:latency}
\end{figure}

\section{Unattended and Service Behaviour}

The service-mode host was installed, the machine rebooted, and a controller
connected before any interactive login --- confirming that presence begins at
boot rather than at logon. Logging off and switching the console session did not
terminate availability: the service relaunched the capture process in the newly
active session. Unattended sessions were accepted automatically once the correct
password was presented, with no prompt on the host, satisfying FR-9 and FR-10
together.

\section{Discussion}

The measurements support the central design claim. Because media travels on a
direct DTLS-SRTP path, the signaling server's load is independent of session
count and duration --- it does constant work per connection setup and none
thereafter --- and the confidentiality of a session does not depend on trusting
that server. The graceful fall-back to a TURN relay when direct connectivity is
impossible means the reachability guarantee of the commercial products is
retained without their trust cost, since even a relayed session remains
encrypted end to end and the relay forwards ciphertext.

Two honest limitations emerged. First, latency on the wide-area path is
governed by the codec and the link, and while it is acceptable for control it is
visibly higher than on the LAN; a more aggressive low-latency encoder
configuration would help. Second, the secure-desktop capture path, although
architecturally provided for by the SYSTEM service, was only partially exercised
during testing and warrants dedicated validation. Both points are carried
forward into Chapter~7.
